% Options for packages loaded elsewhere
\PassOptionsToPackage{unicode}{hyperref}
\PassOptionsToPackage{hyphens}{url}
\documentclass[
]{article}
\usepackage{xcolor}
\usepackage[margin=1in]{geometry}
\usepackage{amsmath,amssymb}
\setcounter{secnumdepth}{-\maxdimen} % remove section numbering
\usepackage{iftex}
\ifPDFTeX
  \usepackage[T1]{fontenc}
  \usepackage[utf8]{inputenc}
  \usepackage{textcomp} % provide euro and other symbols
\else % if luatex or xetex
  \usepackage{unicode-math} % this also loads fontspec
  \defaultfontfeatures{Scale=MatchLowercase}
  \defaultfontfeatures[\rmfamily]{Ligatures=TeX,Scale=1}
\fi
\usepackage{lmodern}
\ifPDFTeX\else
  % xetex/luatex font selection
\fi
% Use upquote if available, for straight quotes in verbatim environments
\IfFileExists{upquote.sty}{\usepackage{upquote}}{}
\IfFileExists{microtype.sty}{% use microtype if available
  \usepackage[]{microtype}
  \UseMicrotypeSet[protrusion]{basicmath} % disable protrusion for tt fonts
}{}
\makeatletter
\@ifundefined{KOMAClassName}{% if non-KOMA class
  \IfFileExists{parskip.sty}{%
    \usepackage{parskip}
  }{% else
    \setlength{\parindent}{0pt}
    \setlength{\parskip}{6pt plus 2pt minus 1pt}}
}{% if KOMA class
  \KOMAoptions{parskip=half}}
\makeatother
\usepackage{color}
\usepackage{fancyvrb}
\newcommand{\VerbBar}{|}
\newcommand{\VERB}{\Verb[commandchars=\\\{\}]}
\DefineVerbatimEnvironment{Highlighting}{Verbatim}{commandchars=\\\{\}}
% Add ',fontsize=\small' for more characters per line
\usepackage{framed}
\definecolor{shadecolor}{RGB}{248,248,248}
\newenvironment{Shaded}{\begin{snugshade}}{\end{snugshade}}
\newcommand{\AlertTok}[1]{\textcolor[rgb]{0.94,0.16,0.16}{#1}}
\newcommand{\AnnotationTok}[1]{\textcolor[rgb]{0.56,0.35,0.01}{\textbf{\textit{#1}}}}
\newcommand{\AttributeTok}[1]{\textcolor[rgb]{0.13,0.29,0.53}{#1}}
\newcommand{\BaseNTok}[1]{\textcolor[rgb]{0.00,0.00,0.81}{#1}}
\newcommand{\BuiltInTok}[1]{#1}
\newcommand{\CharTok}[1]{\textcolor[rgb]{0.31,0.60,0.02}{#1}}
\newcommand{\CommentTok}[1]{\textcolor[rgb]{0.56,0.35,0.01}{\textit{#1}}}
\newcommand{\CommentVarTok}[1]{\textcolor[rgb]{0.56,0.35,0.01}{\textbf{\textit{#1}}}}
\newcommand{\ConstantTok}[1]{\textcolor[rgb]{0.56,0.35,0.01}{#1}}
\newcommand{\ControlFlowTok}[1]{\textcolor[rgb]{0.13,0.29,0.53}{\textbf{#1}}}
\newcommand{\DataTypeTok}[1]{\textcolor[rgb]{0.13,0.29,0.53}{#1}}
\newcommand{\DecValTok}[1]{\textcolor[rgb]{0.00,0.00,0.81}{#1}}
\newcommand{\DocumentationTok}[1]{\textcolor[rgb]{0.56,0.35,0.01}{\textbf{\textit{#1}}}}
\newcommand{\ErrorTok}[1]{\textcolor[rgb]{0.64,0.00,0.00}{\textbf{#1}}}
\newcommand{\ExtensionTok}[1]{#1}
\newcommand{\FloatTok}[1]{\textcolor[rgb]{0.00,0.00,0.81}{#1}}
\newcommand{\FunctionTok}[1]{\textcolor[rgb]{0.13,0.29,0.53}{\textbf{#1}}}
\newcommand{\ImportTok}[1]{#1}
\newcommand{\InformationTok}[1]{\textcolor[rgb]{0.56,0.35,0.01}{\textbf{\textit{#1}}}}
\newcommand{\KeywordTok}[1]{\textcolor[rgb]{0.13,0.29,0.53}{\textbf{#1}}}
\newcommand{\NormalTok}[1]{#1}
\newcommand{\OperatorTok}[1]{\textcolor[rgb]{0.81,0.36,0.00}{\textbf{#1}}}
\newcommand{\OtherTok}[1]{\textcolor[rgb]{0.56,0.35,0.01}{#1}}
\newcommand{\PreprocessorTok}[1]{\textcolor[rgb]{0.56,0.35,0.01}{\textit{#1}}}
\newcommand{\RegionMarkerTok}[1]{#1}
\newcommand{\SpecialCharTok}[1]{\textcolor[rgb]{0.81,0.36,0.00}{\textbf{#1}}}
\newcommand{\SpecialStringTok}[1]{\textcolor[rgb]{0.31,0.60,0.02}{#1}}
\newcommand{\StringTok}[1]{\textcolor[rgb]{0.31,0.60,0.02}{#1}}
\newcommand{\VariableTok}[1]{\textcolor[rgb]{0.00,0.00,0.00}{#1}}
\newcommand{\VerbatimStringTok}[1]{\textcolor[rgb]{0.31,0.60,0.02}{#1}}
\newcommand{\WarningTok}[1]{\textcolor[rgb]{0.56,0.35,0.01}{\textbf{\textit{#1}}}}
\usepackage{graphicx}
\makeatletter
\newsavebox\pandoc@box
\newcommand*\pandocbounded[1]{% scales image to fit in text height/width
  \sbox\pandoc@box{#1}%
  \Gscale@div\@tempa{\textheight}{\dimexpr\ht\pandoc@box+\dp\pandoc@box\relax}%
  \Gscale@div\@tempb{\linewidth}{\wd\pandoc@box}%
  \ifdim\@tempb\p@<\@tempa\p@\let\@tempa\@tempb\fi% select the smaller of both
  \ifdim\@tempa\p@<\p@\scalebox{\@tempa}{\usebox\pandoc@box}%
  \else\usebox{\pandoc@box}%
  \fi%
}
% Set default figure placement to htbp
\def\fps@figure{htbp}
\makeatother
\setlength{\emergencystretch}{3em} % prevent overfull lines
\providecommand{\tightlist}{%
  \setlength{\itemsep}{0pt}\setlength{\parskip}{0pt}}
\usepackage{bookmark}
\IfFileExists{xurl.sty}{\usepackage{xurl}}{} % add URL line breaks if available
\urlstyle{same}
\hypersetup{
  pdftitle={lab01},
  hidelinks,
  pdfcreator={LaTeX via pandoc}}

\title{lab01}
\author{}
\date{\vspace{-2.5em}}

\begin{document}
\maketitle

\#\#Laboratório 01 \#\#\#Importação e Manipulação de Dados em R Benilton
Carvalho, Guilherme Ludwig, Tatiana Benaglia

\#\#\#\#Introdução Um conjunto de dados no formato tidy beneficia o
analista de dados por permitir a manipulação dos mesmos de uma maneira
unificada. De modo similar, métodos estatísticos são habitualmente
implementados para receber dados neste formato. Desta maneira, a
importação e tratamento de dados visando o referido formato reduzirá a
criação de bancos de dados temporários, evitando problemas difíceis de
diagnosticar.

Os conjuntos de dados apresentados correspondem ao número de casos de
tuberculose observados em alguns países, juntamente com seus tamanhos
populacionais.

\#\#\#\#Objetivos Ao fim deste laboratório, você deverá ser capaz de:

Caracterizar um conjunto de dados no formato tidy; Manipular bases de
dados, para responder perguntas específicas, utilizando o arcabouço
implementado via tidyverse, com uso de verbos do pacote dplyr e
apresentação gráfica de resultados utilizando o pacote ggplot2; Notar
que responder tais perguntas com bases de dados que não estejam no
formato tidy requer trabalho adicional;

\#\#\#\#Manipulação de Dados no Formato Tidy 1.Carregue o pacote
tidyverse

\begin{Shaded}
\begin{Highlighting}[]
\FunctionTok{library}\NormalTok{(tidyverse)}
\end{Highlighting}
\end{Shaded}

\begin{verbatim}
## Warning: pacote 'tidyverse' foi compilado no R versão 4.6.1
\end{verbatim}

\begin{verbatim}
## -- Attaching core tidyverse packages ------------------------ tidyverse 2.0.0 --
## v dplyr     1.2.1     v readr     2.2.0
## v forcats   1.0.1     v stringr   1.6.0
## v ggplot2   4.0.3     v tibble    3.3.1
## v lubridate 1.9.5     v tidyr     1.3.2
## v purrr     1.2.2     
## -- Conflicts ------------------------------------------ tidyverse_conflicts() --
## x dplyr::filter() masks stats::filter()
## x dplyr::lag()    masks stats::lag()
## i Use the conflicted package (<http://conflicted.r-lib.org/>) to force all conflicts to become errors
\end{verbatim}

2.Apresente os bancos de dados table1, table2, table3, table4a e
table4b, distribuídos juntamente com o pacote tidyverse. Para cada banco
de dados, descreva textualmente se ele está no formato tidy e justifique
cada uma de suas respostas.

\begin{Shaded}
\begin{Highlighting}[]
\NormalTok{table1}
\end{Highlighting}
\end{Shaded}

\begin{verbatim}
## # A tibble: 6 x 4
##   country      year  cases population
##   <chr>       <dbl>  <dbl>      <dbl>
## 1 Afghanistan  1999    745   19987071
## 2 Afghanistan  2000   2666   20595360
## 3 Brazil       1999  37737  172006362
## 4 Brazil       2000  80488  174504898
## 5 China        1999 212258 1272915272
## 6 China        2000 213766 1280428583
\end{verbatim}

Resposta: Sim, o banco table1 está em formato tidy, pois as colunas
representam variaveis e as linhas contém apenas uma observação para cada
variável

\begin{Shaded}
\begin{Highlighting}[]
\NormalTok{table2}
\end{Highlighting}
\end{Shaded}

\begin{verbatim}
## # A tibble: 12 x 4
##    country      year type            count
##    <chr>       <dbl> <chr>           <dbl>
##  1 Afghanistan  1999 cases             745
##  2 Afghanistan  1999 population   19987071
##  3 Afghanistan  2000 cases            2666
##  4 Afghanistan  2000 population   20595360
##  5 Brazil       1999 cases           37737
##  6 Brazil       1999 population  172006362
##  7 Brazil       2000 cases           80488
##  8 Brazil       2000 population  174504898
##  9 China        1999 cases          212258
## 10 China        1999 population 1272915272
## 11 China        2000 cases          213766
## 12 China        2000 population 1280428583
\end{verbatim}

Resposta: Sim, o banco table2 está em formato tidy, pois as colunas
representam variaveis e as linhas contém apenas uma observação para cada
variável

\begin{Shaded}
\begin{Highlighting}[]
\NormalTok{table3}
\end{Highlighting}
\end{Shaded}

\begin{verbatim}
## # A tibble: 6 x 3
##   country      year rate             
##   <chr>       <dbl> <chr>            
## 1 Afghanistan  1999 745/19987071     
## 2 Afghanistan  2000 2666/20595360    
## 3 Brazil       1999 37737/172006362  
## 4 Brazil       2000 80488/174504898  
## 5 China        1999 212258/1272915272
## 6 China        2000 213766/1280428583
\end{verbatim}

Resposta: Não, o banco table3 não está em formato tidy, pois a coluna
rate esta classificada como e possue duas inforamções em uma única
variavel, embora numericamente ela possa comunicar uma taxa, as
informações contidas na coluna são referentes ao ``cases'' e a
``population''.

\begin{Shaded}
\begin{Highlighting}[]
\NormalTok{table4a}
\end{Highlighting}
\end{Shaded}

\begin{verbatim}
## # A tibble: 3 x 3
##   country     `1999` `2000`
##   <chr>        <dbl>  <dbl>
## 1 Afghanistan    745   2666
## 2 Brazil       37737  80488
## 3 China       212258 213766
\end{verbatim}

Resposta: Não, aqui temos duas colunas cominucando data, o formato longo
(ou tidy) seria formado por uma tabela contendo país, ano e número de
casos

\begin{Shaded}
\begin{Highlighting}[]
\NormalTok{table4b}
\end{Highlighting}
\end{Shaded}

\begin{verbatim}
## # A tibble: 3 x 3
##   country         `1999`     `2000`
##   <chr>            <dbl>      <dbl>
## 1 Afghanistan   19987071   20595360
## 2 Brazil       172006362  174504898
## 3 China       1272915272 1280428583
\end{verbatim}

Resposta: Não, aqui temos duas colunas cominucando data, o formato longo
(ou tidy) seria formado por uma tabela contendo país, ano e número de
casos

3.Utilizando comandos do pacote dplyr, determine a taxa de ocorrência de
tuberculose para cada 10.000 pessoas. Armazene o resultado em um objeto
chamado taxas.

\begin{Shaded}
\begin{Highlighting}[]
\FunctionTok{library}\NormalTok{(dplyr)}

\NormalTok{dados }\OtherTok{\textless{}{-}}\NormalTok{ table1}
\NormalTok{taxas }\OtherTok{\textless{}{-}}\NormalTok{ dados }\SpecialCharTok{\%\textgreater{}\%}
        \FunctionTok{mutate}\NormalTok{(}\AttributeTok{taxa =}\NormalTok{ cases}\SpecialCharTok{/}\NormalTok{population }\SpecialCharTok{*} \DecValTok{10000}\NormalTok{)}
        
\NormalTok{taxas}
\end{Highlighting}
\end{Shaded}

\begin{verbatim}
## # A tibble: 6 x 5
##   country      year  cases population  taxa
##   <chr>       <dbl>  <dbl>      <dbl> <dbl>
## 1 Afghanistan  1999    745   19987071 0.373
## 2 Afghanistan  2000   2666   20595360 1.29 
## 3 Brazil       1999  37737  172006362 2.19 
## 4 Brazil       2000  80488  174504898 4.61 
## 5 China        1999 212258 1272915272 1.67 
## 6 China        2000 213766 1280428583 1.67
\end{verbatim}

4.Apresente, utilizando comandos do pacote dplyr, o número de casos de
tuberculose por ano.

\begin{Shaded}
\begin{Highlighting}[]
\NormalTok{casos\_ano }\OtherTok{\textless{}{-}}\NormalTok{ dados }\SpecialCharTok{\%\textgreater{}\%} 
              \FunctionTok{group\_by}\NormalTok{(year) }\SpecialCharTok{\%\textgreater{}\%} 
                \FunctionTok{summarise}\NormalTok{(}\AttributeTok{casos =} \FunctionTok{sum}\NormalTok{(cases))}
\NormalTok{casos\_ano}
\end{Highlighting}
\end{Shaded}

\begin{verbatim}
## # A tibble: 2 x 2
##    year  casos
##   <dbl>  <dbl>
## 1  1999 250740
## 2  2000 296920
\end{verbatim}

5.Apresente, utilizando comandos do pacote dplyr, o número de casos de
tuberculose identificados em cada país.

\begin{Shaded}
\begin{Highlighting}[]
\NormalTok{casos\_pais }\OtherTok{\textless{}{-}}\NormalTok{ dados }\SpecialCharTok{\%\textgreater{}\%} 
                \FunctionTok{group\_by}\NormalTok{(country) }\SpecialCharTok{\%\textgreater{}\%} 
                  \FunctionTok{summarise}\NormalTok{(}\AttributeTok{total =} \FunctionTok{sum}\NormalTok{(cases))}
\NormalTok{casos\_pais}
\end{Highlighting}
\end{Shaded}

\begin{verbatim}
## # A tibble: 3 x 2
##   country      total
##   <chr>        <dbl>
## 1 Afghanistan   3411
## 2 Brazil      118225
## 3 China       426024
\end{verbatim}

6.Utilizando comandos do pacote dplyr, apresente uma tabela que descreva
a mudança no número de casos, em cada país, ao longo dos anos de 1999 e
2000.

\begin{Shaded}
\begin{Highlighting}[]
\NormalTok{mudanca\_numero\_casos }\OtherTok{\textless{}{-}}\NormalTok{ dados }\SpecialCharTok{\%\textgreater{}\%} 
                          \FunctionTok{arrange}\NormalTok{(country, year) }\SpecialCharTok{\%\textgreater{}\%} 
                            \FunctionTok{group\_by}\NormalTok{(country) }\SpecialCharTok{\%\textgreater{}\%} 
                              \FunctionTok{mutate}\NormalTok{(}\AttributeTok{mudanca =}\NormalTok{ cases }\SpecialCharTok{{-}} \FunctionTok{lag}\NormalTok{(cases)) }\SpecialCharTok{\%\textgreater{}\%} 
                                \FunctionTok{filter}\NormalTok{(year }\SpecialCharTok{==} \DecValTok{2000}\NormalTok{) }\SpecialCharTok{\%\textgreater{}\%} 
                                  \FunctionTok{select}\NormalTok{(country, mudanca)}

\NormalTok{mudanca\_numero\_casos}
\end{Highlighting}
\end{Shaded}

\begin{verbatim}
## # A tibble: 3 x 2
## # Groups:   country [3]
##   country     mudanca
##   <chr>         <dbl>
## 1 Afghanistan    1921
## 2 Brazil        42751
## 3 China          1508
\end{verbatim}

7.Apresente um gráfico de linhas, preparado via ggplot2, apresentando a
mudança na taxa de casos (por 10.000 habitantes) estratificado por país.

\begin{Shaded}
\begin{Highlighting}[]
\FunctionTok{library}\NormalTok{(ggplot2)}
\NormalTok{grafico }\OtherTok{\textless{}{-}} \FunctionTok{ggplot}\NormalTok{(taxas, }\FunctionTok{aes}\NormalTok{(}\AttributeTok{x =}\NormalTok{ year, }\AttributeTok{y =}\NormalTok{ taxa, }\AttributeTok{group =}\NormalTok{ country, }\AttributeTok{colour =}\NormalTok{ country)) }\SpecialCharTok{+}
            \FunctionTok{geom\_line}\NormalTok{() }\SpecialCharTok{+}
              \FunctionTok{geom\_point}\NormalTok{() }\SpecialCharTok{+}
               \FunctionTok{labs}\NormalTok{(}\AttributeTok{x =} \StringTok{"Ano"}\NormalTok{, }\AttributeTok{y =} \StringTok{"Taxa de casos por 10 000 habitantes"}\NormalTok{, }\AttributeTok{color =} \StringTok{"País"}\NormalTok{) }\SpecialCharTok{+}
                  \FunctionTok{theme\_minimal}\NormalTok{()}
\NormalTok{grafico}
\end{Highlighting}
\end{Shaded}

\pandocbounded{\includegraphics[keepaspectratio]{lab01_files/lab01_files/figure-latex/unnamed-chunk-11-1.pdf}}
8. Calcule a taxa para as tabelas table2 e table4a+table4b. Para isso,
você precisará executar 4 passos:

Extrair o número de casos de tuberculose por país, por ano; Extrair o
tamanho da população correspondente, por ano; Dividir o número de casos
pelo tamanho da população e multiplicar o resultado por 10.000;
Armazenar o resultado numa variável apropriada;

\begin{Shaded}
\begin{Highlighting}[]
\CommentTok{\#para table2}

\NormalTok{casos\_table2 }\OtherTok{\textless{}{-}}\NormalTok{ table2 }\SpecialCharTok{\%\textgreater{}\%} 
                  \FunctionTok{filter}\NormalTok{(type }\SpecialCharTok{==} \StringTok{"cases"}\NormalTok{)}

\NormalTok{populacao\_table2 }\OtherTok{\textless{}{-}}\NormalTok{ table2 }\SpecialCharTok{\%\textgreater{}\%} 
                      \FunctionTok{filter}\NormalTok{(type }\SpecialCharTok{==} \StringTok{"population"}\NormalTok{)}
\NormalTok{taxa\_tabela2 }\OtherTok{\textless{}{-}}\NormalTok{ casos\_table2 }\SpecialCharTok{\%\textgreater{}\%} 
                  \FunctionTok{mutate}\NormalTok{(}\AttributeTok{taxa =}\NormalTok{ count}\SpecialCharTok{/}\NormalTok{populacao\_table2}\SpecialCharTok{$}\NormalTok{count }\SpecialCharTok{*} \DecValTok{10000}\NormalTok{)}
\NormalTok{taxa\_tabela2}
\end{Highlighting}
\end{Shaded}

\begin{verbatim}
## # A tibble: 6 x 5
##   country      year type   count  taxa
##   <chr>       <dbl> <chr>  <dbl> <dbl>
## 1 Afghanistan  1999 cases    745 0.373
## 2 Afghanistan  2000 cases   2666 1.29 
## 3 Brazil       1999 cases  37737 2.19 
## 4 Brazil       2000 cases  80488 4.61 
## 5 China        1999 cases 212258 1.67 
## 6 China        2000 cases 213766 1.67
\end{verbatim}

\begin{Shaded}
\begin{Highlighting}[]
\CommentTok{\#para table4a + table4b}

\NormalTok{casos\_table4a }\OtherTok{\textless{}{-}}\NormalTok{ table4a }\SpecialCharTok{\%\textgreater{}\%}
                  \FunctionTok{pivot\_longer}\NormalTok{(}
                     \AttributeTok{cols =} \FunctionTok{c}\NormalTok{(}\StringTok{\textasciigrave{}}\AttributeTok{1999}\StringTok{\textasciigrave{}}\NormalTok{, }\StringTok{\textasciigrave{}}\AttributeTok{2000}\StringTok{\textasciigrave{}}\NormalTok{),}
                     \AttributeTok{names\_to =} \StringTok{"year"}\NormalTok{,}
                     \AttributeTok{values\_to =} \StringTok{"cases"}
\NormalTok{  )}

\NormalTok{casos\_table4a}
\end{Highlighting}
\end{Shaded}

\begin{verbatim}
## # A tibble: 6 x 3
##   country     year   cases
##   <chr>       <chr>  <dbl>
## 1 Afghanistan 1999     745
## 2 Afghanistan 2000    2666
## 3 Brazil      1999   37737
## 4 Brazil      2000   80488
## 5 China       1999  212258
## 6 China       2000  213766
\end{verbatim}

\begin{Shaded}
\begin{Highlighting}[]
\NormalTok{populacao\_table4b }\OtherTok{\textless{}{-}}\NormalTok{ table4b }\SpecialCharTok{\%\textgreater{}\%}
                      \FunctionTok{pivot\_longer}\NormalTok{(}
                        \AttributeTok{cols =} \FunctionTok{c}\NormalTok{(}\StringTok{\textasciigrave{}}\AttributeTok{1999}\StringTok{\textasciigrave{}}\NormalTok{, }\StringTok{\textasciigrave{}}\AttributeTok{2000}\StringTok{\textasciigrave{}}\NormalTok{),}
                        \AttributeTok{names\_to =} \StringTok{"year"}\NormalTok{,}
                        \AttributeTok{values\_to =} \StringTok{"population"}
\NormalTok{  )}

\NormalTok{populacao\_table4b}
\end{Highlighting}
\end{Shaded}

\begin{verbatim}
## # A tibble: 6 x 3
##   country     year  population
##   <chr>       <chr>      <dbl>
## 1 Afghanistan 1999    19987071
## 2 Afghanistan 2000    20595360
## 3 Brazil      1999   172006362
## 4 Brazil      2000   174504898
## 5 China       1999  1272915272
## 6 China       2000  1280428583
\end{verbatim}

\begin{Shaded}
\begin{Highlighting}[]
\NormalTok{taxa\_tabela4 }\OtherTok{\textless{}{-}}\NormalTok{ casos\_table4a }\SpecialCharTok{\%\textgreater{}\%} 
                  \FunctionTok{mutate}\NormalTok{(}\AttributeTok{taxa =}\NormalTok{ cases}\SpecialCharTok{/}\NormalTok{ populacao\_table4b}\SpecialCharTok{$}\NormalTok{population }\SpecialCharTok{*} \DecValTok{10000}\NormalTok{)}
\NormalTok{taxa\_tabela4}
\end{Highlighting}
\end{Shaded}

\begin{verbatim}
## # A tibble: 6 x 4
##   country     year   cases  taxa
##   <chr>       <chr>  <dbl> <dbl>
## 1 Afghanistan 1999     745 0.373
## 2 Afghanistan 2000    2666 1.29 
## 3 Brazil      1999   37737 2.19 
## 4 Brazil      2000   80488 4.61 
## 5 China       1999  212258 1.67 
## 6 China       2000  213766 1.67
\end{verbatim}

9.Refaça o gráfico da questão 7 para os dados apresentados em table2.

\begin{Shaded}
\begin{Highlighting}[]
\NormalTok{grafico2 }\OtherTok{\textless{}{-}} \FunctionTok{ggplot}\NormalTok{(taxa\_tabela2, }\FunctionTok{aes}\NormalTok{(}\AttributeTok{x =}\NormalTok{ year, }\AttributeTok{y =}\NormalTok{ taxa, }\AttributeTok{group =}\NormalTok{ country, }\AttributeTok{colour =}\NormalTok{ country)) }\SpecialCharTok{+}
            \FunctionTok{geom\_line}\NormalTok{() }\SpecialCharTok{+}
              \FunctionTok{geom\_point}\NormalTok{() }\SpecialCharTok{+}
               \FunctionTok{labs}\NormalTok{(}\AttributeTok{x =} \StringTok{"Ano"}\NormalTok{, }\AttributeTok{y =} \StringTok{"Taxa de casos por 10 000 habitantes"}\NormalTok{, }\AttributeTok{color =} \StringTok{"País"}\NormalTok{) }\SpecialCharTok{+}
                  \FunctionTok{theme\_minimal}\NormalTok{()}
\NormalTok{grafico2}
\end{Highlighting}
\end{Shaded}

\pandocbounded{\includegraphics[keepaspectratio]{lab01_files/lab01_files/figure-latex/unnamed-chunk-14-1.pdf}}

\begin{enumerate}
\def\labelenumi{\arabic{enumi}.}
\setcounter{enumi}{9}
\tightlist
\item
  Utilizando o comando pivot\_longer, transforme table4a em um objeto no
  formato tidy. Armazene o resultado num objeto chamado tidy4a.
\end{enumerate}

\begin{Shaded}
\begin{Highlighting}[]
\NormalTok{tidy4a }\OtherTok{\textless{}{-}}\NormalTok{ table4a }\SpecialCharTok{\%\textgreater{}\%} 
                  \FunctionTok{pivot\_longer}\NormalTok{(}
                    \AttributeTok{cols =} \FunctionTok{c}\NormalTok{(}\StringTok{\textquotesingle{}1999\textquotesingle{}}\NormalTok{, }\StringTok{\textquotesingle{}2000\textquotesingle{}}\NormalTok{),}
                    \AttributeTok{names\_to =} \StringTok{"ano"}\NormalTok{,}
                    \AttributeTok{values\_to =} \StringTok{"casos"}
\NormalTok{                  )}
\end{Highlighting}
\end{Shaded}

11.Refaça o item 10 para o objeto table4b. Armazene o resultado num
objeto chamado tidy4b.

\begin{Shaded}
\begin{Highlighting}[]
\NormalTok{tidy4b }\OtherTok{\textless{}{-}}\NormalTok{ table4b }\SpecialCharTok{\%\textgreater{}\%} 
                  \FunctionTok{pivot\_longer}\NormalTok{(}
                    \AttributeTok{cols =} \FunctionTok{c}\NormalTok{(}\StringTok{\textquotesingle{}1999\textquotesingle{}}\NormalTok{,}\StringTok{\textquotesingle{}2000\textquotesingle{}}\NormalTok{),}
                    \AttributeTok{names\_to =} \StringTok{"ano"}\NormalTok{,}
                    \AttributeTok{values\_to =} \StringTok{"população"}
\NormalTok{                  )}
\end{Highlighting}
\end{Shaded}

12.Combine os objetos tidy4a e tidy4b em um único objeto, utilizando o
comando left\_join. Apresente uma explicação textual sobre o que faz o
referido comando.

\begin{Shaded}
\begin{Highlighting}[]
\NormalTok{tabela4 }\OtherTok{\textless{}{-}} \FunctionTok{left\_join}\NormalTok{(}
\NormalTok{            tidy4a,}
\NormalTok{            tidy4b,}
            \AttributeTok{by =} \FunctionTok{c}\NormalTok{(}\StringTok{"country"}\NormalTok{, }\StringTok{"ano"}\NormalTok{)}
\NormalTok{)}

\NormalTok{tabela4}
\end{Highlighting}
\end{Shaded}

\begin{verbatim}
## # A tibble: 6 x 4
##   country     ano    casos  população
##   <chr>       <chr>  <dbl>      <dbl>
## 1 Afghanistan 1999     745   19987071
## 2 Afghanistan 2000    2666   20595360
## 3 Brazil      1999   37737  172006362
## 4 Brazil      2000   80488  174504898
## 5 China       1999  212258 1272915272
## 6 China       2000  213766 1280428583
\end{verbatim}

Resposta: o comando left\_join compara as linhas da tabela à esquerda
com as da direita, mantém as linhas da esquerda e adiciona as linhas
diferentes da tabela à direita, formando assim uma única tabela sem
conflitos de informações. Caso não haja correspondência, ele não cria
uma observação e sim adiciona uma celula NA.No exemplo acima as colunas
comparadas foram ``country'' e ``ano'', ou seja, quando essas duas
informações coincidiram ele pode completar a linha de observações unindo
as inforamções contidas nas duas tabelas.

13.Use o comando pivot\_wider para tranformar o objeto table2 em um
objeto com formato tidy.

\begin{Shaded}
\begin{Highlighting}[]
\NormalTok{tabela2 }\OtherTok{\textless{}{-}}\NormalTok{ table2 }\SpecialCharTok{\%\textgreater{}\%} 
            \FunctionTok{pivot\_wider}\NormalTok{(}
              \AttributeTok{names\_from =}\NormalTok{ type,}
              \AttributeTok{values\_from =}\NormalTok{ count}
\NormalTok{            )}

\NormalTok{tabela2}
\end{Highlighting}
\end{Shaded}

\begin{verbatim}
## # A tibble: 6 x 4
##   country      year  cases population
##   <chr>       <dbl>  <dbl>      <dbl>
## 1 Afghanistan  1999    745   19987071
## 2 Afghanistan  2000   2666   20595360
## 3 Brazil       1999  37737  172006362
## 4 Brazil       2000  80488  174504898
## 5 China        1999 212258 1272915272
## 6 China        2000 213766 1280428583
\end{verbatim}

14.Observe que a coluna rate do objeto table3 é um texto mostrando a
fração que formaria a taxa de casos de tuberculose. Transforme o objeto
table3 em um objeto com formato tidy separando a coluna 3 em duas outras
colunas: cases e population, utilizando o comando separate. Utilize o
argumento convert para transformar o resultado em um objeto numérico.

\begin{Shaded}
\begin{Highlighting}[]
\NormalTok{tabela3 }\OtherTok{\textless{}{-}}\NormalTok{ table3 }\SpecialCharTok{\%\textgreater{}\%} 
            \FunctionTok{separate}\NormalTok{(}
\NormalTok{              rate,}
              \AttributeTok{into =} \FunctionTok{c}\NormalTok{(}\StringTok{"cases"}\NormalTok{, }\StringTok{"population"}\NormalTok{),}
              \AttributeTok{sep =} \StringTok{"/"}\NormalTok{,}
              \AttributeTok{convert =} \ConstantTok{TRUE}
\NormalTok{            )}
\NormalTok{tabela3}
\end{Highlighting}
\end{Shaded}

\begin{verbatim}
## # A tibble: 6 x 4
##   country      year  cases population
##   <chr>       <dbl>  <int>      <int>
## 1 Afghanistan  1999    745   19987071
## 2 Afghanistan  2000   2666   20595360
## 3 Brazil       1999  37737  172006362
## 4 Brazil       2000  80488  174504898
## 5 China        1999 212258 1272915272
## 6 China        2000 213766 1280428583
\end{verbatim}

\end{document}
